% {\color{blue} 초안 작성 v0.0}

% {\color{blue} subsection에 대한 고민 필요}

% {\color{blue} A.비교군 / B.평가지표설명 / C.RL구조 / D.Sched\_Ext로 구성. }

\section{Fair Scheduling and Virtual Runtime}

Linux의 Completely Fair Scheduler(CFS)는 각 Task의 Scheduling
Weight에 비례하여 CPU 자원을 분배하는 것을 목적으로 한다.
CFS는 실제 CPU Execution Time을 Scheduling Weight에 따라
정규화한 Virtual Runtime(Vruntime)을 이용하여 각 Task의 상대적인
CPU 사용량을 관리한다.
Task $i$가 $\Delta t_i$만큼 CPU를 사용하였을 때 Vruntime의
증가량은 다음과 같이 표현할 수 있다.

\begin{equation}
    V_i(t+\Delta t_i)-V_i(t) = \Delta t_i \frac{w_0}{w_i},
\end{equation}

여기서 $w_i$는 Task $i$의 Scheduling Weight이며,
$w_0$는 nice 0에 해당하는 기준 Scheduling Weight이다.
높은 Scheduling Weight를 가진 Task는 동일한 CPU Execution Time에 대해
Vruntime이 상대적으로 적게 증가하므로 더 많은 CPU 실행 기회를
확보할 수 있다.
CFS는 동일 Runqueue 내에서 Vruntime이 작은 Scheduling Entity를
우선적으로 선택함으로써 장기적인 CPU 자원 분배의 Fairness를
유지한다.

Linux 6.6부터 도입된 Earliest Eligible Virtual Deadline
First(EEVDF)는 기존 Vruntime에 기반한 CPU 사용량 관리에
Eligibility와 Virtual Deadline이라는 개념을 추가한다.
EEVDF는 할당받아야 할 CPU 자원보다 적은 자원을 사용한
Eligible Task 중 가장 이른 Virtual Deadline을 가진 Task를
우선적으로 선택한다.
그러나 Virtual Deadline은 Scheduling Weight와 요청된
CPU 실행 구간에 기반하여 계산되는 가상 시간상의 값으로,
Application Request의 실제 Deadline과는 구분된다.
따라서 CFS와 EEVDF에서 정의하는 실행 우선순위는
Application Request의 시간적 제약과 반드시 일치하지 않는다.


\section{Request-Level Deadlines and Task Dependencies}

Task는 CPU Scheduler가 관리하는 실행 단위이며,
Request는 Application이 완료해야 하는 작업 단위이다.
Workload는 이러한 Task와 Request 및 이들의 실행 패턴과
의존 관계를 포함하는 전체 작업 구성을 의미한다.
하나의 Task는 여러 Request를 반복적으로 처리할 수 있으며,
반대로 하나의 Request가 완료되기 위해 여러 Task의 실행과
Synchronization이 필요할 수도 있다.
예를 들어 Multimedia Application의 Frame Processing은
여러 Task의 연산, Wakeup, Blocking 및 Synchronization을
포함할 수 있다.

Request $j$의 Arrival Time, Completion Time 및
Absolute Deadline을 각각 $a_j$, $c_j$, $d_j$라고 정의한다.
Deadline Budget은 Request의 Arrival Time부터 Absolute Deadline까지
허용된 전체 시간이며, Remaining Deadline은 현재 시각 $t$에서
Absolute Deadline까지 남은 시간을 의미한다.

\begin{equation}
    \mathrm{Deadline\ Budget}_j = d_j-a_j,\qquad
    \mathrm{Remaining\ Deadline}_j(t) = d_j-t.
    \label{eq:deadline_times}
\end{equation}

Deadline Budget은 Request마다 정해지는 값인 반면,
Remaining Deadline은 시간에 따라 감소하며 Deadline이 지나면 음수가 된다.
이와 달리 Remaining CPU Demand는 Request를 완료하기 위해
추가로 필요한 CPU 실행 시간을 의미한다.

Request Completion Latency($L_j$), Deadline Miss($M_j$),
Deadline Overrun($O_j$)은 다음과 같이 표현할 수 있다.

\begin{equation}
    L_j = c_j - a_j,\qquad M_j = \mathbf{1}[c_j > d_j],\qquad O_j = \max(0, c_j - d_j).
\end{equation}

Deadline Miss는 Request가 시간적 제약을 충족하지 못하였는지를
나타내며, Deadline Overrun은 해당 제약을 얼마나 초과하였는지를
나타낸다.
Deadline Miss만을 최소화하는 정책은 이미 Deadline을 초과한
Request의 실행을 지연시키더라도 추가적인 불이익을 받지 않을 수
있으므로, Request Completion Latency와 Deadline Overrun을
함께 고려할 필요가 있다.

본 논문에서 Scheduling Urgency는 Request의 Deadline Compliance를
향상시키기 위해 해당 Task를 우선 실행할 필요성을 의미한다.
Scheduling Priority는 Scheduling Policy가 결정하는 상대적인 실행
우선순위를 의미하므로 Scheduling Urgency와 구분된다.
Request의 Deadline과 해당 Request에 관여하는 Task의
Scheduling Urgency는 반드시 일치하지 않는다.
예를 들어 Remaining Deadline이 짧은 Request를 처리하는 Task가
Lock을 기다리고 있다면, 해당 Task를 직접 실행하는 대신
Lock Holder를 먼저 실행하는 것이 Request Completion을
앞당길 수 있다.
따라서 Task의 Scheduling Urgency는 개별 Request의
Remaining Deadline뿐만 아니라 Remaining CPU Demand, Task 간 Dependency 및
현재의 CPU Contention에 의해서도 결정된다.


\section{Extensible Scheduling with sched\_ext}

Linux 6.12에 도입된 sched\_ext는 eBPF Scheduling Program을 통해
사용자 정의 CPU Scheduling Policy를 구현할 수 있도록
지원하는 Scheduler Class이다.
sched\_ext는 CPU Selection, Task Enqueue 및 Dispatch 등의
Scheduling Operation에 대한 Interface를 제공한다.

sched\_ext는 Dispatch Queue(DSQ)를 이용하여 eBPF Scheduling Program과
Kernel Scheduler Core 사이에서 Runnable Task를 관리한다.
DSQ는 FIFO 또는 Virtual Time 기반의 Priority Queue로
동작할 수 있으며, 이를 통해 사용자 정의 Scheduling Metric에
따른 Task Ordering을 구현할 수 있다.
본 논문에서 Task Ordering은 Task 간 상대적인 실행 순서를,
Task Selection은 다음 실행할 Task를 선택하는 행위를,
Task Dispatch는 선택된 Task를 Kernel의 실제 실행 경로에
반영하는 과정을 의미한다. Scheduling Decision은 이러한 실행 대상에
대한 Scheduling Policy의 의사결정을 의미한다.

그러나 CPU Scheduling은 빈번하게 수행되는 Kernel Critical
Path에 해당하므로, Scheduling Policy의 추가적인 연산 비용은
전체 시스템의 성능에 직접적인 영향을 미칠 수 있다.
또한 eBPF Scheduling Program은 Verifier를 통해 메모리 접근의 안전성과
프로그램 실행의 유효성을 검증받아야 하므로, 복잡한 연산을
포함하는 Machine Learning Model을 구현하는 데 제약이 존재한다.

따라서 학습 기반 Scheduling Policy를 Kernel에 적용하기
위해서는 Model의 표현력뿐만 아니라 Inference Cost와
Scheduling Decision의 실행 비용을 함께 고려해야 한다.


\section{Learning-based Scheduling Techniques}

\paragraph{Reinforcement Learning (RL).}
RL은 Agent가 Environment와 상호작용하면서 State에 따른
Action을 결정하고, 그 결과로 주어지는 Reward를 통해
Policy를 학습하는 방법이다.
RL의 목적은 개별 Action의 즉각적인 Reward뿐만 아니라
이후의 누적 Reward를 최대화하는 Policy를 학습하는 것이다.
CPU Scheduling에서는 현재의 Task Selection이 이후 Task의
Runnable Waiting Time, Wakeup 및 Request Completion에 연쇄적으로 영향을
미치므로, RL을 통해 연속적인 Scheduling Decision의
장기적인 성능 영향을 고려할 수 있다.

\paragraph{DeepSets.}
DeepSets는 원소의 개수가 가변적이고 입력 순서에 의미가
없는 집합 형태의 데이터를 처리하기 위한 Neural Network
Architecture이다.
각 원소를 공유된 Encoder $\phi$를 통해 개별적으로
Encoding하고, 이를 순서에 독립적인 Aggregation
Operation을 통해 하나의 Context로 표현한다.

\begin{equation}
    z = \rho \left( \sum_{i=1}^{N} \phi(x_i) \right),
\end{equation}

여기서 $x_i$는 집합의 개별 원소이며, $\rho$는
Aggregation된 Feature를 변환하는 함수이다.
이러한 구조는 입력 원소의 개수와 순서가 달라지더라도
집합 전체의 특성을 고정된 차원의 Representation으로
표현할 수 있다.
CPU Scheduling에서는 경쟁하는 Runnable Task의 수가
동적으로 변화하므로, DeepSets를 통해 개별 Task의
특성과 경쟁 Task들의 Context를 함께 고려할 수 있다.

\paragraph{Teacher--Student Knowledge Distillation (KD).}
Teacher--Student Knowledge Distillation은 높은 표현력을 가진 Teacher
Model의 학습된 Knowledge를 상대적으로 경량화된
Student Model에 전달하는 방법이다.
Student는 Teacher의 Output 또는 Policy Distribution을
학습 목표로 사용하여 Teacher의 동작을 근사한다.
이를 통해 Training 단계에서는 높은 연산 비용을 요구하는
복잡한 Model을 활용하고, Deployment 단계에서는 제한된
연산 자원에 적합한 경량 Model을 실행할 수 있다.
이러한 Training과 Inference의 분리는 실행 비용에
민감한 CPU Scheduling 환경에서 학습 기반 Policy를
적용하는 데 유용하다.
